\documentclass[12pt,a4paper]{article}
\usepackage{fontspec}
\setmainfont{Kinnari}[Path=C:/project/thesis-latex/fonts/, Extension=.otf, BoldFont=*-Bold, ItalicFont=*-Italic, BoldItalicFont=*-BoldItalic, Script=Thai]
\setmonofont{texgyretermes}[Path=C:/project/thesis-latex/fonts/, Extension=.otf, UprightFont=*-regular]
\XeTeXlinebreaklocale "th_TH"
\XeTeXlinebreakskip = 0pt plus 1pt
\usepackage[margin=2.2cm]{geometry}
\usepackage{booktabs,array,float,enumitem,xcolor,tcolorbox}
\usepackage[hidelinks]{hyperref}
\usepackage{fancyhdr,titlesec}
\titleformat{\section}{\Large\bfseries}{\thesection}{0.7em}{}
\titleformat{\subsection}{\large\bfseries}{\thesubsection}{0.7em}{}
\pagestyle{fancy}
\fancyhf{}
\fancyhead[L]{\small สรุปผลการประเมินระบบ Agentic AI ด้านการเงินส่วนบุคคล}
\fancyhead[R]{\small 7 พฤศจิกายน 2569}
\fancyfoot[C]{\small \thepage}
\setlength{\parskip}{0.5em}
\setlength{\parindent}{0pt}
\renewcommand{\arraystretch}{1.25}
\setlist{itemsep=0.2em, topsep=0.3em}

\newcolumntype{L}[1]{>{\raggedright\arraybackslash}p{#1}}
\definecolor{boxblue}{HTML}{EEF3FA}
\definecolor{boxline}{HTML}{3B6EA8}
\newtcolorbox{takeaway}{colback=boxblue, colframe=boxline, boxrule=0.6pt, arc=2pt, left=6pt, right=6pt, top=4pt, bottom=4pt, fontupper=\normalsize}

\begin{document}

\begin{center}
{\LARGE\bfseries สรุปผลการประเมินระบบ\par}
\vspace{0.3em}
{\large Agentic AI ด้านการเงินส่วนบุคคลสำหรับผู้ใช้ชาวไทย\par}
\vspace{0.3em}
{\normalsize Pakkaphon (รหัสนักศึกษา 66070146) \quad คณะเทคโนโลยีสารสนเทศ สถาบันเทคโนโลยีพระจอมเกล้าเจ้าคุณทหารลาดกระบัง\par}
\end{center}

\section*{สรุปผลใน 1 หน้า}

ระบบเป็นสถาปัตยกรรมแบบหลาย Agent (Multi-Agent Architecture) ที่มี Router Agent จำแนกเจตนาของคำถาม แล้วส่งต่อให้ Agent เฉพาะทาง 6 ตัว หรือ General Chat ระบบที่ใช้งานจริงใช้ Google Gemini 3.5 Flash คณะผู้จัดทำประเมินระบบ 4 ด้าน ผลโดยสรุปเป็นดังตาราง

\begin{table}[H]
\centering
\small
\begin{tabular}{L{2.6cm}L{5.2cm}L{7.6cm}}
\toprule
\textbf{การทดสอบ} & \textbf{คำถามที่ต้องการตอบ} & \textbf{ผลลัพธ์} \\
\midrule
สถาปัตยกรรม & แบบใดทำงานสำเร็จมากที่สุด & บน DeepSeek แบบ Single Agent สำเร็จ 98.9\% ระบบจริง (แบบ Hub-and-Spoke ที่ Multi-intent Routing) สำเร็จ 92.2\% (แบบไม่แยกงาน 87.2\%) โดยใช้ token น้อยกว่าแบบ Single Agent 3.3 เท่า \\
Multi-intent Routing & รองรับคำถามที่ต้องใช้หลาย Agent ได้หรือไม่ & Router ที่ใช้งานจริงจำแนกถูก 205 จาก 215 ข้อ (95.3\%) ทั้งคำถามเรื่องเดียวและหลายเรื่อง คำถามที่ต้องใช้หลาย Agent สำเร็จประมาณ 86.1\% เทียบกับ 61.1\% ของแบบไม่แยกงาน ($p = 0.012$) \\
ระดับการคิดของ Router & ใช้ระดับต่ำสุดแล้วยังแม่นยำหรือไม่ & แม่นยำไม่ต่างกัน (95.9\% เทียบกับ 96.2\%) และเร็วขึ้นเกือบเท่าตัว (1.36 เทียบกับ 2.59 วินาที) \\
การค้นคืนเอกสาร & โมเดลเวกเตอร์แทนความหมายตัวใดดีที่สุด & Gemini ค้นคืนได้ดีที่สุด (nDCG@10 0.805) แต่คำตอบสุดท้ายของ 3 โมเดลที่ดีที่สุดถูกต้องใกล้เคียงกัน (92.5--94.2\%) \\
การคำนวณภาษี & ภาษีถูกต้องตามกฎปี 2569 หรือไม่ & ถูกต้อง 26 จาก 26 กรณี ถ้าไม่ใช้เครื่องมือคำนวณ โมเดลตอบถูก 7 จาก 8 กรณี และยังใช้เพดานประกันสังคมแบบเก่า \\
\bottomrule
\end{tabular}
\end{table}

\begin{takeaway}
\textbf{ข้อสรุปหลัก} \quad แบบ Hub-and-Spoke ที่มี Multi-intent Routing ปิดช่องว่างกับแบบ Single Agent ได้เกือบหมด โดยยังประหยัด token สำหรับคำถามเรื่องเดียว ส่วนการตั้งค่า Router การเลือกโมเดลเวกเตอร์ และการคำนวณภาษี ผลการทดลองรองรับการออกแบบที่ใช้อยู่ทั้งหมด
\end{takeaway}

ตัวเลขทุกตัวในรายงานนี้มาจากไฟล์ผลการทดลองโดยตรง และทุกรอบที่ล้มเหลวในการทดสอบสถาปัตยกรรมได้อ่านตรวจด้วยมือก่อนนับผล ข้อมูลรายข้อทั้งหมดอยู่ในรายงานฉบับเต็มภาษาอังกฤษ (\texttt{evaluation-report.pdf})

\clearpage

\section{วิธีทดสอบโดยย่อ}

ระบบสร้างโมเดลภาษาผ่านส่วนสร้างโมเดลกลางที่รองรับ Google Gemini, Ollama และ OpenRouter จึงเปลี่ยนโมเดลได้ด้วยการตั้งค่าเพียงจุดเดียว การทดสอบสถาปัตยกรรมใช้โมเดล 2 ตระกูล คือ gpt-oss:20b และ DeepSeek V4.1 Flash ส่วนการทดสอบอื่นใช้ Gemini ตัวเดียวกับระบบจริง

\begin{table}[H]
\centering
\small
\begin{tabular}{lrL{9.5cm}}
\toprule
\textbf{ชุดข้อมูล} & \textbf{จำนวน} & \textbf{รายละเอียด} \\
\midrule
งานทดสอบสถาปัตยกรรม & 60 & คำถามเรื่องเดียว 40 งาน ใช้หลาย Agent 12 งาน และคำถามกว้างที่ต้องถามกลับ 8 งาน \\
คำถามจำแนกเจตนา & 185 & เขียนใหม่ 150 ข้อ และชุดเดือนกันยายน 35 ข้อ \\
ข้อความที่ถามหลายเรื่อง & 30 & หลายเรื่อง 22 ข้อ (2--5 เรื่อง) และคำถามเดียวที่ฟังคล้ายหลายเรื่อง 8 ข้อ \\
คำถามค้นคืนเอกสาร & 120 & ใช้ภาษาแบบผู้ใช้ พร้อมไฟล์และคำสำคัญที่ใช้ตัดสินว่าเอกสารเกี่ยวข้อง \\
กรณีภาษี & 26 & กรณีทั่วไป 20 กรณี และกรณีตามกฎปี 2569 อีก 6 กรณี \\
\bottomrule
\end{tabular}
\end{table}

การเปรียบเทียบใช้ช่วงความเชื่อมั่น 95\% แบบ Wilson และใช้การทดสอบ McNemar แบบแม่นตรงเมื่อเทียบงานเดียวกันในรอบเดียวกัน ค่า pass$^3$ คือสัดส่วนงานที่สำเร็จครบทั้ง 3 รอบ ซึ่งใช้วัดความสม่ำเสมอ เวลาที่รายงานทั้งหมดเป็นเวลาตอบโดยทั่วไป (มัธยฐาน) คือค่าที่อยู่ตรงกลางเมื่อเรียงเวลาของทุกรอบจากน้อยไปมาก ซึ่งไม่ถูกดึงขึ้นจากรอบที่ช้าผิดปกติเหมือนค่าเฉลี่ย

\section{การเปรียบเทียบสถาปัตยกรรม}

คณะผู้จัดทำสร้างสถาปัตยกรรม 4 แบบจาก Agent และเครื่องมือชุดเดียวกัน ได้แก่ แบบ Single Agent ที่มีเครื่องมือทั้ง 27 รายการ แบบ Hub-and-Spoke ที่ใช้ในระบบจริง แบบ Hierarchical และแบบ Peer-to-Peer ทุกแบบทำงาน 60 งาน งานละ 3 รอบ งานหนึ่งนับว่าสำเร็จเมื่อส่งถูก Agent เรียกเครื่องมือที่จำเป็นครบ ตอบตัวเลขถูก และบันทึกข้อมูลถูกต้อง

\begin{table}[H]
\centering
\small
\begin{tabular}{L{4.6cm}rrrrr}
\toprule
\textbf{สถาปัตยกรรม (DeepSeek)} & \textbf{สำเร็จรวม} & \textbf{ใช้หลาย Agent} & \textbf{pass$^3$} & \textbf{Token ขาเข้า} & \textbf{\begin{tabular}[b]{@{}r@{}}เวลาตอบทั่วไป\\(วินาที)\end{tabular}} \\
\midrule
Single Agent & 98.9\% & 100.0\% & 96.7\% & 42,555 & 3.6 \\
Hub-and-Spoke + Multi-intent Routing (ระบบจริง)$^{*}$ & 92.2\% & 86.1\% & 90.0\% & 12,946 & 4.4 \\
แบบ Peer-to-Peer & 91.7\% & 66.7\% & 88.3\% & 14,807 & 5.6 \\
Hub-and-Spoke (ไม่แยกงาน) & 87.2\% & 61.1\% & 83.3\% & 11,649 & 4.4 \\
แบบ Hierarchical & 85.0\% & 55.6\% & 80.0\% & 12,280 & 5.3 \\
\bottomrule
\end{tabular}
\end{table}

$^{*}$ แถวระบบจริงวัด Multi-intent Routing จาก 6 งานที่ต้องใช้หลาย Agent ซึ่งแบบไม่แยกงานล้มเหลวอย่างน้อย 1 รอบ (18 รอบ) ส่วนงานอื่นผ่านเส้นทางเดียวกันจึงใช้ผลร่วมกัน (ดูหัวข้อ Multi-intent Routing) ระบบจริงยังต่ำกว่าแบบ Single Agent อย่างมีนัยสำคัญ (14 คู่เทียบกับ 2 คู่, $p = 0.004$) แต่ pass$^3$ อยู่ที่ 90.0\% เทียบกับ 83.3\% ของแบบไม่แยกงาน และเวลาเท่ากัน

บน gpt-oss:20b (ไม่ได้ทดสอบ Multi-intent Routing) ผลต่างออกไป แบบ Single Agent และแบบ Hub-and-Spoke เสมอกัน (83.9\% และ 83.3\%, $p = 1.0$) และแบบ Hub-and-Spoke สม่ำเสมอที่สุด (pass$^3$ 78.3\% เทียบกับ 73.3\%)

\begin{takeaway}
\textbf{ข้อสรุป} \quad กับโมเดลที่เก่ง แบบ Single Agent ชนะเพราะใช้เครื่องมือข้ามด้านได้ในคำถามเดียว ส่วนแบบ Hub-and-Spoke ที่ไม่แยกงานส่งคำถามให้ Agent ตัวเดียว จึงตอบคำถามที่ต้องใช้หลาย Agentไม่ครบ ข้อดีของแบบ Hub-and-Spoke คือใช้ token น้อยกว่า 3.6 เท่า และไม่เสียเปรียบเมื่อใช้โมเดลที่อ่อนกว่า
\end{takeaway}

\section{Multi-intent Routing}

Multi-intent Routing ให้ Router Agent แยกข้อความที่ถามหลายเรื่องเป็นงานย่อย งานละ 1 Agent ตามลำดับที่ผู้ใช้ถาม โดยไม่จำกัดจำนวนงาน Agent ทำงานทีละตัว ตัวหลังเห็นคำตอบของตัวก่อนหน้า แล้วระบบรวมคำตอบเป็นข้อความเดียว คำถามเรื่องเดียวส่งให้ Agent ตัวเดียวตามปกติ

กลไกนี้อยู่ใน Router Agent ทั้งหมด คณะผู้จัดทำจึงทดสอบดังนี้

\begin{table}[H]
\centering
\small
\begin{tabular}{L{4.3cm}L{11cm}}
\toprule
\textbf{การทดสอบ} & \textbf{ผลลัพธ์} \\
\midrule
Router ที่ใช้งานจริง 215 ข้อ & ถูก 205/215 = 95.3\% (CI 91.7--97.5\%): คำถามเรื่องเดียว 175/185 (94.6\%, ไม่ต่างจาก 95.7\% ของคำสั่งที่ไม่มี Multi-intent Routing, $p = 0.69$, แยกเกิน 1 ข้อ) หลายเรื่อง 22/22 ถูกทุก Agent และลำดับ และคำถามเดียวที่ฟังคล้ายหลายเรื่อง 8/8 \\
ทำงานจริงบน DeepSeek & ทดสอบ 6 งานที่ต้องใช้หลาย Agent ซึ่งแบบไม่แยกงานล้มเหลวอย่างน้อย 1 รอบ (18 รอบ) ส่วนงานอื่นใช้ผลร่วมกัน \\
\bottomrule
\end{tabular}
\end{table}

\begin{table}[H]
\centering
\small
\begin{tabular}{L{6.4cm}rrr}
\toprule
\textbf{สถาปัตยกรรม (DeepSeek)} & \textbf{สำเร็จรวม} & \textbf{ใช้หลาย Agent} & \textbf{ใช้หลาย Agent pass$^3$} \\
\midrule
Hub-and-Spoke (ไม่แยกงาน) & 87.2\% & 61.1\% & 6/12 \\
แบบ Hub-and-Spoke + Multi-intent Routing (ประมาณการ) & 92.2\% & 86.1\% & 10/12 \\
Single Agent & 98.9\% & 100.0\% & 12/12 \\
\bottomrule
\end{tabular}
\end{table}

ใน 18 คู่นี้ Multi-intent Routing สำเร็จในจุดที่แบบไม่แยกงานล้มเหลว 10 ครั้ง และกลับกันเพียง 1 ครั้ง ($p = 0.012$) ข้อแลกเปลี่ยนคือข้อความที่ถามหลายเรื่องใช้ทรัพยากรมากขึ้น โดยเรียกโมเดลเพิ่มจาก 3.4 เป็น 5.8 ครั้ง token จาก 11.2k เป็น 24.1k และเวลาจาก 9.0 เป็น 13.0 วินาที ส่วนคำถามเรื่องเดียวไม่ได้รับผลกระทบ

\begin{takeaway}
\textbf{ข้อควรระวัง} \quad ตัวเลข 86.1\% เป็นค่าประมาณ เพราะไม่ได้ทดสอบ 6 งานที่ต้องใช้หลาย Agent ซึ่งแบบไม่แยกงานสำเร็จทุกรอบ ความล้มเหลวที่เหลือ 5 รอบอ่านตรวจด้วยมือแล้ว คือส่งคำถามเรื่องการออมไปผิด Agent 2 รอบ และ Agent ภาษีของ DeepSeek ตอบว่าง 3 รอบ
\end{takeaway}

\section{ระดับการคิดของ Router Agent}

บทที่ 3 ระบุว่า Router Agent ใช้ระดับการคิดต่ำสุดของ Gemini โดยไม่เสียความแม่นยำ การทดสอบนี้ใช้คำถาม 185 ข้อ ข้อละ 2 รอบ ผลการทดสอบเบื้องต้นพบว่าระดับ low และ minimal ใช้ token สำหรับการคิดเป็นศูนย์เท่ากัน จึงเทียบ minimal กับค่าเริ่มต้นของโมเดล

\begin{table}[H]
\centering
\small
\begin{tabular}{llrrr}
\toprule
\textbf{โมเดล} & \textbf{ระดับการคิด} & \textbf{ความแม่นยำ} & \textbf{เวลาตอบทั่วไป (วินาที)} & \textbf{Token การคิด} \\
\midrule
gemini-3.5-flash & minimal (ระบบจริง) & 95.9\% & 1.36 & 0 \\
gemini-3.5-flash & ไม่ตั้งค่า & 96.2\% & 2.59 & 272 \\
gemini-3.5-flash-lite & minimal & 93.8\% & 0.88 & 0 \\
gemini-3.5-flash-lite & high & 96.5\% & 1.93 & 469 \\
\bottomrule
\end{tabular}
\end{table}

\begin{takeaway}
\textbf{ข้อสรุป} \quad ระดับ minimal ผ่านเกณฑ์ไม่ด้อยกว่า ความแม่นยำต่างกัน $-0.3$ จุด ช่วงความเชื่อมั่นอยู่ที่ $-2.2$ ถึง $+1.6$ จุด ซึ่งอยู่เหนือเกณฑ์ $-3$ จุดทั้งช่วง และเร็วขึ้นประมาณ 1.26 วินาที รองรับข้อความในบทที่ 3 ส่วน Flash-Lite แม่นยำน้อยกว่า 2.2 จุด จึงไม่นำมาใช้ ส่วนคำสั่งของ Router ที่ใช้งานจริงซึ่งมี Multi-intent Routing จำแนกคำถามเรื่องเดียวถูก 94.6\% ไม่ต่างจากผลในตารางนี้อย่างมีนัยสำคัญ
\end{takeaway}

\section{การค้นคืนเอกสาร}

คณะผู้จัดทำเทียบวิธีค้นคืน 5 แบบ ด้วยคำถาม 120 ข้อ บนฐานความรู้ที่ตัดข้อความแบบเดียวกับระบบจริง จากนั้นให้ Gemini ตอบคำถามจากเอกสาร 3 ชิ้นแรกของแต่ละวิธี และให้ DeepSeek ซึ่งเป็นโมเดลต่างตระกูลเป็นผู้ตัดสิน ผู้ตัดสินให้ผลตรงกับการอ่านตรวจด้วยมือ 28 จาก 30 คำตอบ

\begin{table}[H]
\centering
\small
\begin{tabular}{lrrrr}
\toprule
\textbf{วิธีค้นคืน} & \textbf{nDCG@10} & \textbf{Recall@3} & \textbf{คำตอบถูกต้อง} & \textbf{ตรงกับแหล่งอ้างอิง} \\
\midrule
gemini-embedding-001 (ระบบจริง) & 0.805 & 0.900 & 93.3\% & 88.3\% \\
BAAI/bge-m3 & 0.724 & 0.842 & 94.2\% & 82.5\% \\
multilingual-e5-large-instruct & 0.722 & 0.825 & 92.5\% & 86.7\% \\
BM25 & 0.568 & 0.683 & -- & -- \\
MiniLM & 0.394 & 0.433 & -- & -- \\
\bottomrule
\end{tabular}
\end{table}

\begin{takeaway}
\textbf{ข้อสรุป} \quad Gemini ค้นคืนได้ดีกว่าทุกวิธีอย่างมีนัยสำคัญ จึงใช้ต่อในระบบจริง ส่วน MiniLM ไม่เหมาะกับภาษาไทย ระบบจึงใช้ bge-m3 เป็นตัวสำรองบนเครื่อง อย่างไรก็ตาม ความถูกต้องของคำตอบสุดท้ายของ 3 วิธีแรกไม่ต่างกันอย่างมีนัยสำคัญ เพราะโมเดลมักตอบถูกได้แม้เอกสาร 3 ชิ้นจะเกี่ยวข้องเพียงบางส่วน
\end{takeaway}

\section{การคำนวณภาษี}

ตัวคำนวณภาษีและคำสั่งของ Tax Agent ใช้กฎปี 2569 ชุดทดสอบจึงมี 6 กรณีที่ถ้าใช้กฎของปีก่อนจะได้คำตอบผิด

\begin{table}[H]
\centering
\small
\begin{tabular}{L{4cm}rL{8.6cm}}
\toprule
\textbf{การทดสอบ} & \textbf{ผล} & \textbf{หมายเหตุ} \\
\midrule
ความถูกต้องผ่าน Tax Agent & 26/26 & ไม่มีข้อผิดพลาด ช่วงความเชื่อมั่นของอัตราผิดพลาดอยู่ที่ 0--12.9\% \\
ใช้เครื่องมือคำนวณ & 8/8 & ถูกต้องทุกกรณี \\
ไม่ใช้เครื่องมือ (อ่านด้วยมือ) & 7/8 & ผิดกรณีเพดานรวม RMF กับ SSF (ตอบ 467,000 แทน 407,000 บาท) และยังใช้เพดานประกันสังคม 9,000 บาทแบบเก่า \\
การตรวจข้อมูลผิด 10 กรณี & 10/10 & ระบบตรวจอัตโนมัติให้ 9/10 แต่เมื่ออ่านแล้วกรณีที่ไม่ผ่านตอบถูกต้อง \\
\bottomrule
\end{tabular}
\end{table}

\begin{takeaway}
\textbf{ข้อสรุป} \quad ยังควรใช้เครื่องมือคำนวณภาษี ไม่ใช่เพราะโมเดลคิดเลขไม่ได้ แต่เพราะเครื่องมือทำให้ทุกคำตอบถูกต้องและเป็นกฎปัจจุบันเสมอ
\end{takeaway}

\section{ข้อจำกัด}

\begin{itemize}
\item ไม่ได้ทดสอบสถาปัตยกรรมซ้ำบน Gemini 3.5 Flash ด้วยข้อจำกัดด้านงบประมาณ คณะผู้จัดทำเลือก Gemini สำหรับระบบจริงเพราะเสถียรที่สุดในการทดสอบที่ใช้ Gemini คือ ภาษีถูก 26 จาก 26 กรณี จำแนกคำถาม 185 ข้อโดยไม่มีครั้งใดล้มเหลว และไม่พบคำตอบว่าง ขณะที่ DeepSeek ตอบว่างเป็นระยะ และ gpt-oss:20b เคยถามกลับเป็นภาษาจีน
\item ผลของ Multi-intent Routing ในคำถามที่ต้องใช้หลาย Agentเป็นค่าประมาณ
\item ชุดข้อมูลทั้งหมดเขียนขึ้นเอง ไม่ได้มาจากข้อความของผู้ใช้จริง
\item ยังไม่มีการประเมินโดยผู้เชี่ยวชาญ การทดสอบกับผู้ใช้จริง หรือการทดสอบเชิงโจมตี
\item เวลาตอบวัดจากเครื่องของคณะผู้จัดทำ ไม่ได้วัดจากเซิร์ฟเวอร์จริง
\end{itemize}

\end{document}
